% !PRINTED 228
\weberpage{228}
\begin{center}
\S\,67. \textbf{\gesperrt{Besondere Transformationsgleichungen.}}
\end{center}

Es sollen nun die Wurzeln der Transformationsgleichungen durch $\vartheta$-Funktionen dargestellt werden. Wir setzen zu diesem Zweck:

\weq{1}{
\begin{aligned}
\Omega &= \frac{4\mu K + 4\mu' i K'}{n}, \quad iK' = K\omega,\\
\varpi &= \frac{\mu+\mu'\omega}{n}, \quad \Omega = 4K\varpi
\end{aligned}
}

und betrachten die Produkte

\[
\prod_{1,\,\frac{n-1}{2}}^{h} \vartheta_{00}(2h\varpi), \quad
\prod_{1,\,\frac{n-1}{2}}^{h} \vartheta_{01}(2h\varpi), \quad
\prod_{1,\,\frac{n-1}{2}}^{h} \vartheta_{10}(2h\varpi), \quad
\prod_{1,\,\frac{n-1}{2}}^{h} \vartheta_{11}(2h\varpi).
\]

Es empfiehlt sich folgende Bezeichnung:

\weq{2}{
\begin{aligned}
P_{00}\,\vartheta_{00}^{\frac{n-1}{2}} &= e^{\frac{\pi i}{6}\mu'(\mu+\mu'\omega)\frac{n^2-1}{n}} \prod_{1,\,\frac{n-1}{2}}^{h} \vartheta_{00}(2h\varpi),\\
P_{10}\,\vartheta_{10}^{\frac{n-1}{2}} &= e^{\frac{\pi i}{6}\mu'(\mu+\mu'\omega)\frac{n^2-1}{n}} \prod_{1,\,\frac{n-1}{2}}^{h} \vartheta_{10}(2h\varpi),\\
P_{01}\,\vartheta_{01}^{\frac{n-1}{2}} &= e^{\frac{\pi i}{6}\mu'(\mu+\mu'\omega)\frac{n^2-1}{n}} \prod_{1,\,\frac{n-1}{2}}^{h} \vartheta_{01}(2h\varpi).
\end{aligned}
}

Es ist jetzt die Formel \S\,39, (9) anzuwenden, die man aber für den gegenwärtigen Zweck etwas anders darstellt.

Setzt man in \S\,39, (8) $v' = 2h/n$, $v = 2h(\alpha+\beta\omega)/n$ und nimmt abermals das Produkt über $h = 1, 2, \ldots (n-1)/2$, so folgt mit Anwendung von \S\,32, (24):

\[
\begin{aligned}
&2^{\frac{n-1}{2}} \prod_{1,\,\frac{n-1}{2}}^{h} \vartheta_{00}\!\left(\frac{2h(\alpha+\beta\omega)}{n}\right) \vartheta_{01}\!\left(\frac{2h(\alpha+\beta\omega)}{n}\right) \vartheta_{10}\!\left(\frac{2h(\alpha+\beta\omega)}{n}\right)\\
&= (-1)^{\frac{n^2-1}{8}} e^{-\frac{\pi i}{2}\frac{n^2-1}{n}\beta(\alpha+\beta\omega)} \vartheta_{00}^{\frac{n-1}{2}} \vartheta_{10}^{\frac{n-1}{2}} \vartheta_{01}^{\frac{n-1}{2}}\textsuperscript{1)}.
\end{aligned}
\]

Demnach erhält man durch Multiplikation der drei Gleichungen (2):

\weq{3}{(-1)^{\frac{n^2-1}{8}} 2^{\frac{n-1}{2}} P_{00} P_{10} P_{01} = 1.}

\textsuperscript{1)}~In \S\,34, (9) der ersten Auflage ist die betreffende Formel gleich in dieser Form angegeben, die sich auch, wenn auch etwas umständlicher, aus \S\,39, (9) ableiten läßt.
